\subsection{Taylor's formula for formal power series}

Let \(\mathbf{X} = (\mathbf{X}_i)_{i\in I}\) and \(\mathbf{Y} = (\mathbf{Y}_i)_{i\in I}\) be two families of indeterminates relative to the same index set I. We denote by \(\mathbf{X} + \mathbf{Y}\) the family \((\mathbf{X}_i + \mathbf{Y}_i)_{i\in I}\) of formal power series in \(\mathbf{A}[[\mathbf{X},\mathbf{Y}]]\) . It is clear that we can substitute \(\mathbf{X}_i + \mathbf{Y}_i\) for \(\mathbf{X}_i\) in a formal power series \(u\in \mathbf{A}[[\mathbf{X}]]\) , the result being written \(u(\mathbf{X} + \mathbf{Y})\) . For each \(\nu \in \mathbf{N}^{(I)}\) we denote by \(\Delta^{\nu}u\) the coefficient of \(\mathbf{Y}^{\nu}\) in the formal power series \(u(\mathbf{X} + \mathbf{Y})\) considered as belonging to \(\mathbf{A}[[\mathbf{X}]][[\mathbf{Y}]]\) (III, p.~456). In other words, we have

\[
u (\mathbf {X} + \mathbf {Y}) = \sum_ {\nu} \Delta^ {\nu} u (\mathbf {X}) \cdot \mathbf {Y} ^ {\nu} \quad (u \in \mathrm{A} [ [ \mathbf {X} ] ]).\tag{6}
\]

Substituting (0, X) for (X, Y) we obtain

\[
u (\mathbf {X}) = \sum_ {\nu} \Delta^ {\nu} u (0) \cdot \mathbf {X} ^ {\nu};\tag{7}
\]

In other words, the constant term of \(\Delta^{\nu}u\) is the coefficient of \(\mathbf{X}^{\nu}\) in \(u\) . Since the mapping \(u \mapsto u(\mathbf{X} + \mathbf{Y})\) of \(\mathbf{A}[[\mathbf{X}]]\) into \(\mathbf{A}[[\mathbf{X}, \mathbf{Y}]]\) is continuous, the mappings \(u \mapsto \Delta^{\nu}u\) of \(\mathbf{A}[[\mathbf{X}]]\) into itself are again continuous.

As in the case of polynomials (IV, p.~7) we can prove the formulae

\begin{enumerate}
\def\labelenumi{(\arabic{enumi})}
\setcounter{enumi}{7}
\tightlist
\item
\end{enumerate}

\[
\Delta^ {\sigma} (u v) = \sum_ {\nu + \rho = \sigma} \Delta^ {\nu} (u) \Delta^ {\rho} (v),
\]

\[
\Delta^ {\rho} \Delta^ {\sigma} u = \frac {(\rho + \sigma) !}{\rho ! \sigma !} \Delta^ {\rho + \sigma} u.\tag{9}
\]

The binomial formula (I, p.~99, Cor. 2) gives the following value for \(\Delta^{\nu}u\) when \(u = \sum_{\lambda} \alpha_{\lambda} X^{\lambda}\)

\[
\Delta^ {\nu} u = \sum_ {\lambda} \alpha_ {\lambda + \nu} \frac {(\lambda + \nu) !}{\lambda ! \nu !} X ^ {\lambda}.\tag{10}
\]

Consider in particular the case \(\nu = \varepsilon_{i}\) , that is \(\nu_{i} = 1\) , \(\nu_{j} = 0\) for \(j \neq i\) . We shall put \(D_{i}u = \Delta^{\varepsilon_{i}}u\) ; put differently, \(D_{i}u\) is the coefficient of \(Y_{i}\) in \(u(\mathbf{X} + \mathbf{Y})\) . By (10) we thus have

\[
\mathrm{D} _ {i} u = \sum_ {\lambda} (\lambda_ {i} + 1) \alpha_ {\lambda + \varepsilon_ {i}} X ^ {\lambda};\tag{11}
\]

in particular we have \(\mathbf{D}_{i}(\mathbf{X}_{i})=1\) and \(\mathbf{D}_{i}(\mathbf{X}_{j})=0\) for \(j\neq i\) . The formula (8) shows that \(D_{i}\) is a derivation of A{[}{[}X{]}{]}, and from (9) we deduce the relation

\[
\mathbf {D} ^ {\nu} \boldsymbol {u} = \nu ! \Delta^ {\nu} \boldsymbol {u}\tag{12}
\]

as in the case of polynomials (IV, p.~8) (we have put \(\mathbf{D}^{\nu} = \prod_{i\in I}\mathbf{D}_{i}^{\nu_{i}}\) for \(\nu = (\nu_{i})_{i\in I}\) in \(\mathbf{N}^{(I)}\) ). When \(\mathbf{A}\) is a \(\mathbf{Q}\) -algebra, the formulae (6), (7) and (12) imply the « Taylor formulae »:

\begin{enumerate}
\def\labelenumi{(\arabic{enumi})}
\setcounter{enumi}{12}
\tightlist
\item
\end{enumerate}

\[
u (\mathbf {X} + \mathbf {Y}) = \sum_ {\nu} \frac {1}{\nu !} D ^ {\nu} u (\mathbf {X}). \mathbf {Y} ^ {\nu},
\]

\[
u (\mathbf {X}) = \sum_ {\nu} \frac {1}{\nu !} D ^ {\nu} u (0) \cdot X ^ {\nu}.\tag{14}
\]

Remarks. --- 1) We often say that \(D_i u\) is the partial derivative of \(u\) with respect to \(X_i\) ; we also use the notation \(D_{X_i} u, \frac{\partial u}{\partial X_i}\) and \(u_{X_i}'\) . For a single indeterminate \(X\) the unique partial derivative \(Du\) (also written \(\frac{du}{dX}\) or \(u'\) ) is called the derivative of \(u\) .

\begin{enumerate}
\def\labelenumi{\arabic{enumi})}
\setcounter{enumi}{1}
\item
  The formula (9) shows that the endomorphisms \(\Delta^{\rho}\) of the A-module \(\mathbf{A}[[\mathbf{X}]]\) commute pairwise. Hence the same holds of the endomorphisms \(D_i\) .
\item
  If \(u \in \mathbf{A}[(\mathbf{X}_i)_{i \in I}]\) is a polynomial, the polynomials \(\Delta^{\rho} u\) and \(D_i u\) defined in IV, p.~6 and 7 coincide with the formal power series denoted by the same symbols.
\end{enumerate}

